\section{相关工作}
\subsection{稀疏专家模型与路由}
Switch Transformer 以稀疏路由将每个 token 分配给少量专家，从而在控制单次计算量的同时扩大模型容量\cite{fedus2022switch}。GShard 进一步将稀疏专家计算与自动分片结合，说明专家模型的可扩展性不仅取决于算子计算量，也取决于跨设备分片和通信组织\cite{lepikhin2021gshard}。这类工作主要关注模型训练、路由均衡或分布式执行；本文关注的问题更窄，即在固定的 3D CIM 空间中，将已知或可获得的专家激活轨迹转化为权重驻留、复制和静态调度决策。

DeepSeek-V3 技术报告展示了大规模 MoE 模型中路由专家、共享专家和负载平衡机制的组合形式\cite{deepseek2024v3}。该结构为本文的建模提供了实际背景，但本文不复现 DeepSeek-V3 的完整权重、训练过程或推理栈；项目中的 128 专家实验只用于验证 metadata-only 工程链路的可运行性。

\subsection{CIM 映射与加速器评估}
CIM 和近存计算的核心动机是减少权重与计算之间的数据移动。已有加速器研究系统总结了数据流、片上存储、并行阵列和数据复用对 DNN 加速器性能的影响\cite{sze2017efficient}。Timeloop 则提供了面向 DNN 加速器的数据流和映射空间评估思路\cite{parashar2019timeloop}。这些工作为本文的“资源映射—调度模拟—指标分析”方法提供了方法论参照，但并未直接给出本赛题规定的 $N^2$ Sub-Cube、Weight-Cube 分层放置和同一 Sub-Cube 单一激活约束。

因此，本文的区别在于把三类信息放进同一个可验证闭环：一是由模型形状得到的几何切分，二是由激活轨迹得到的专家关系统计，三是由静态任务图得到的依赖和资源占用。该闭环的目标不是提出通用 CIM 芯片架构，而是为赛题二提供一个输入、映射、调度、验证和结果固化均可追踪的工程方案。
